%-------------------------------------------------------------------------------
%   CHAPTER 1 - УДИРТГАЛ
%-------------------------------------------------------------------------------

\section{Удиртгал}
\label{sec:introduction}

\subsection{Судалгааны үндэслэл}
\label{subsec:background}

Монгол бичиг нь Монгол үндэстний түүх, соёлын үнэт өв болох уламжлалт босоо бичгийн систем юм. XIII зууны эхэн үед бүрэлдэн тогтож, төрийн албан бичиг болсноос хойш энэхүү бичгийн систем нь олон зууны турш Монголчуудын соёл, шашин, гүн ухаан, шинжлэх ухааны мэдлэгийг тээж, уламжлан дамжуулсаар ирсэн~\cite{poppe1974mongolian}. Гэвч 1941 онд Монгол улс албан ёсоор кирилл үсгийг төрийн бичиг болгон баталснаар Монгол бичгийн өдөр тутмын хэрэглээ эрс хумигдаж, орчин үеийн нийгэмд, ялангуяа залуу үеийнхний дунд уламжлалт бичгээ эзэмшсэн иргэдийн тоо тасралтгүй буурсаар байна~\cite{janhunen2003mongolic}.

\paragraph{Монгол бичгийн өнөөгийн нөхцөл байдал}

Монгол улсын Засгийн газраас 2020 онд ``Монгол бичгийг нэвтрүүлэх үндэсний хөтөлбөр''-ийг батлан гаргажээ. Гэвч одоогийн байдлаар орчин үеийн мэдээллийн технологид тулгуурласан дижитал сургалтын хэрэглэгдэхүүн нэн хомс байна. Уламжлалт заах арга зүй нь үр дүнтэй биш бөгөөд залуу үеийнхний анхаарлыг татахгүй байгаа юм.

\paragraph{Хэл сурах технологийн эрчимтэй хөгжил}
Дэлхий дахинд гадаад хэл болон төрөлх хэлээ гүнзгийрүүлэн суралцах зориулалттай ухаалаг гар утасны аппликейшнууд (Duolingo, Memrise, Anki зэрэг) өргөнөөр нэвтэрч, хэл эзэмших үйл явцыг илүү хүртээмжтэй, үр дүнтэй бөгөөд сонирхолтой хэлбэрт шилжүүлсээр байна~\cite{godwin2016mobile}. Ялангуяа, ``Spaced repetition'' буюу зайтай давталтын аргачлалд суурилсан Flashcard (санах ойн карт) систем нь тархины мэдээлэл хадгалах зүй тогтолд тулгуурладаг тул шинэ үг хэллэг цээжлэх, танин мэдэхүйн чадварыг сайжруулахад шинжлэх ухааны хувьд өндөр ач холбогдолтой болох нь олон улсын судалгаануудаар батлагдсан~\cite{nakata2011computer}. Хэдийгээр технологийн ийм дэвшил гарсаар байгаа боловч, Монгол бичгийг бие даан суралцахад зориулагдсан орчин үеийн, цогц арга зүй бүхий аппликейшн одоогоор зах зээлд үгүйлэгдэж байна.

\paragraph{Шийдвэрлэх шаардлагатай үндсэн асуудлууд}

Монгол бичгийг түгээн дэлгэрүүлэхэд дараах үндсэн хүндрэлүүд тулгарч байна:
\begin{enumerate}
    \item \textbf{Сургалтын материалын дутагдал:} Орчин үеийн интерактив сургалтын хэрэглэгдэхүүн дутагдалтай байна.
    \item \textbf{Мотивацийн дутагдал:} Зөвхөн цээжлэхэд суурилсан уламжлалт аргачлал нь сэдлийг хөндөхгүй байна.
\end{enumerate}

Эдгээр тулгамдсан асуудлууд нь Монгол бичгийг эзэмших болон өдөр тутамдаа хэрэглэх иргэдийн тоог бууруулах, улмаар үндэсний бичиг соёлын дархлааг сулруулах эрсдэлийг дагуулж байна. Тиймээс орчин үеийн гар утасны технологи болон ухаалаг алгоритмд суурилсан, хэрэглэгч төвтэй дизайн бүхий интерактив Монгол бичиг сургалтын системийг хөгжүүлэх нь тулгамдсан бөгөөд цаг үеэ олсон зайлшгүй шаардлага болж байна.

\paragraph{Одоо байгаа системүүдийн дутагдал}

Дэлхий даяар хэл сурах аппликейшнууд (Duolingo, Anki, Memrise) өргөн тархсан боловч Монгол бичгийг таних, гар бичлэгийг анализ хийх, соёлын контекст олгох функцүүд дутагдалтай байна.

\paragraph{Egüü системийн давуу талууд}

Egüü систем нь Монгол бичгийн орчин үеийн апп бөгөөд: (1) TensorFlow OCR-ээр гар бичлэг таних, (2) SM-2 Spaced Repetition + Gamification, (3) Орчин үеийн UI/UX, (4) Соёлын контекст, (5) Үнэгүй, нээлттэй эх зэрэг онцлогтай.

\subsection{Судалгааны зорилго}
\label{subsec:objective}

Энэхүү төгсөлтийн судалгааны ажлын үндсэн зорилго нь ухаалаг төхөөрөмжийн дэвшилтэт технологиудыг (ялангуяа машин сургалт болон гар утасны аппликейшн хөгжүүлэлтийн сүүлийн үеийн чиг хандлагыг) ашиглан Монгол бичгийг сонирхолтой, үр дүнтэйгээр бие даан суралцах боломжийг олгох \textbf{``Egüü''} (босоо Монгол бичгийн ``үсэг'' хэмээх утгаар) хэмээх гар утасны аппликейшныг системийн түвшинд зохион бүтээж, хөгжүүлэхэд оршино.

\noindent Энэхүү систем нь дараах үндсэн функцуудыг хэрэгжүүлнэ:
\begin{itemize}
    \item Flashcard системээр Монгол бичгийн үсгүүдийг зураг, фонетик, орчуулгын хамт суралцах
    \item TensorFlow OCR ашиглан хэрэглэгчийн гараар бичсэн Монгол бичгийг таних
    \item Өдөр бүрийн шинэ үг, streak tracking системээр мотивацийг нэмэгдүүлэх
    \item Олон хэлний дэмжлэг (Монгол, Англи), dark/light theme зэрэг орчин үеийн UI/UX
    \item Firebase/Firestore backend ашиглан хэрэглэгчийн өгөгдлийг хадгалах, синхрончлох
\end{itemize}

\subsection{Судалгааны зорилтууд}
\label{subsec:tasks}

Ерөнхий зорилгыг амжилттай хэрэгжүүлэхийн тулд дараах үе шаттай нарийвчилсан зорилтуудыг дэвшүүлж, шийдвэрлэнэ:

\begin{enumerate}
    \item \textbf{Онолын болон харьцуулсан макро түвшний судалгаа хийх:}
    \begin{itemize}
        \item Монгол бичгийн үүсэл гарвал, хэл зүйн бүтэц болон заах арга зүйн онцлогийг судлах.
        \item Flashcard систем болон "Spaced repetition" алгоритмын танин мэдэхүйн онолын үндсийг судлах.
        \item Одоо зах зээлд тэргүүлж буй хэл сурах аппликейшнуудын (Duolingo, Anki, Memrise) давуу болон сул талуудад харьцуулсан дүн шинжилгээ хийх.
        \item Системд ашиглагдах React Native, Firebase, TensorFlow OCR зэрэг технологиудын чадавхи, хязгаарлалтыг үнэлэх.
    \end{itemize}
    
    \item \textbf{Системийн архитектурын шинжилгээ ба зохиомж боловсруулах:}
    \begin{itemize}
        \item Программ хангамжийн инженерингийн аргачлалаар хэрэглэгчийн шаардлагыг тодорхойлж, UML (Unified Modeling Language) диаграммуудыг боловсруулах.
        \item Системийн ерөнхий архитектур болон үүлэн өгөгдлийн сангийн (Cloud Firestore) зохиомжийн шийдлийг гаргах.
        \item Орчин үеийн чиг хандлагад нийцсэн хэрэглэгчийн интерфэйс (UI/UX) болон хэрэглэгчийн туршлагын загварыг боловсруулах.
        \item Сургалтад ашиглагдах Flashcard өгөгдлийн сан буюу dataset-ийг системтэйгээр бэлтгэх.
    \end{itemize}
    
    \item \textbf{Цогц системийн програм хангамжийн хэрэгжүүлэлт:}
    \begin{itemize}
        \item React Native (Expo фреймворк) ашиглан iOS болон Android үйлдлийн системүүдэд зэрэг ажиллах (cross-platform) аппликейшн хөгжүүлэх.
        \item Firebase (Authentication, Firestore, Cloud Functions) үйлчилгээнүүдэд түшиглэн аюулгүй, өргөжүүлэх боломжтой backend дэд бүтцийг угсрах.
        \item TensorFlow.js ашиглан ухаалаг төхөөрөмж дээр шууд ажиллах Монгол бичгийн үсэг таних OCR машины сургалтын загварыг нэгтгэх.
        \item NativeWind (Tailwind CSS) ашиглан респонсив UI хөгжүүлэх болон i18n сангаар олон хэлний орчны харилцан үйлчлэлийг хангах.
    \end{itemize}
    
    \item \textbf{Нарийвчилсан туршилт болон хүн-компьютерийн харилцан үйлчлэлийн үнэлгээ:}
    \begin{itemize}
        \item Системийн найдвартай ажиллагааг хангах зорилгоор модулиудын үйл ажиллагааг техникийн хувьд тестлэх.
        \item Бодит хэрэглэгчдийн дунд туршилт зохион байгуулж, хэрэглэгчийн туршлага, интерфейсийн ойлгомжтой байдлыг үнэлэх.
        \item Хэрэглэгчийн сэтгэл ханамжийн судалгаа авч, системийн давуу тал болон цаашид сайжруулах шаардлагатайг тодорхойлох.
    \end{itemize}
\end{enumerate}

\subsection{Судалгааны ач холбогдол}
\label{subsec:significance}

Энэхүү судалгааны ажлын үр дүнд хөгжүүлсэн ``Egüü'' гар утасны аппликейшн нь програм хангамжийн инженерчлэл болон сурган хүмүүжүүлэх ухааны огтлолцолд орших бөгөөд дараах практик, нийгэм соёлын өндөр ач холбогдлыг агуулж байна:

\subsubsection{Практик болон шинжлэх ухааны ач холбогдол}
\begin{itemize}
    \item \textbf{Технологид суурилсан боловсролын түгээлт:} Зайтай давталтын алгоритм (Spaced repetition) болон бичил сургалтын (Microlearning) аргуудыг нэвтрүүлснээр суралцагчийн мэдлэг эзэмших хурд болон тогтоох чадварыг мэдэгдэхүйц нэмэгдүүлнэ.
    
    \item \textbf{Хүртээмжтэй байдлыг хангах:} Газарзүйн болон цаг хугацааны хязгаарлалтгүйгээр, ухаалаг утастай хэн бүхэнд Монгол бичгээр бие даан суралцах дижитал орчинг бүрдүүлнэ.
    
    \item \textbf{Дэвшилтэт технологийн интеграци:} Гар утасны хэрэглээнд зориулагдсан хөнгөн жингийн машин сургалтын (TensorFlow.js) загварыг Монгол бичиг танихад ашигласан нь хиймэл оюун ухааныг боловсролын салбарт нэвтрүүлэх амжилттай жишиг проектын нэг болно.
    
    \item \textbf{Суралцах тасралтгүй үйл явцыг дэмжих:} Дараалсан өдрийн тооцоолол (Streak tracking), өдрийн үг (Daily word) зэрэг тоглоомчлолын (Gamification) элементүүд нь хэрэглэгчийн мотивацийг өдөөж, суралцах дадлыг төлөвшүүлнэ.
\end{itemize}

\subsubsection{Нийгэм, соёлын ач холбогдол}
\begin{itemize}
    \item \textbf{Үндэсний өвийг дижитал хэлбэрээр өвлүүлэх:} Монгол үндэстний үнэлж баршгүй соёл, түүхийн өв болох босоо бичгийг орчин үеийн залууст хамгийн ойр хэрэглэгдэхүүн болох гар утсаар дамжуулан заах нь уг өвийг хадгалан үлдэх, сэргээн хөгжүүлэх онцгой ач холбогдолтой.
    
    \item \textbf{Төрийн бодлогыг дэмжих:} ``Монгол бичгийг нэвтрүүлэх үндэсний хөтөлбөр''-ийн хүрээнд дэвшүүлсэн зорилтуудыг хэрэгжүүлэхэд шаардлагатай цахим дэд бүтэц, хэрэглэгдэхүүний санг баяжуулна.
    
    \item \textbf{Боловсролын салбарын цахим шилжилт:} Мэдээлэл, харилцаа холбооны технологийн ололтыг хэл шинжлэл, заах арга зүйд нэвтрүүлэн дижитал шилжилтийн үйл явцыг хурдасгах бодит туршлага болно.
    
    \item \textbf{Олон улсад сурталчлах:} Энэхүү систем нь Монгол хэл, бичгийг сонирхон судалж буй гадаад улс орны иргэдэд хүртээмжтэй, ойлгомжтой сургалтын хэрэглэгдэхүүн болох бөгөөд Монгол хэлийг олон улсад сурталчлан таниулах гүүр болно.
\end{itemize}

\subsection{Судалгааны хамрах хүрээ ба хязгаарлалт}
\label{subsec:scope}

\subsubsection{Хамрах хүрээ}
\begin{itemize}
    \item \textbf{Агуулгын хүрээ:} Монгол бичгийн үндсэн үсгүүд, түгээмэл хэрэглэгддэг үгс (500+ flashcard).
    \item \textbf{Техникийн хүрээ:} iOS болон Android үйлдлийн системтэй ухаалаг гар утас хэрэглэгчид.
    \item \textbf{Хэлний хүрээ:} Монгол, Англи хэлний дэмжлэг.
    \item \textbf{Функцийн хүрээ:} Flashcard систем, OCR таних, streak tracking, daily word, favorites, profile, settings.
\end{itemize}

\subsubsection{Хязгаарлалт}
\begin{itemize}
    \item \textbf{Сүлжээний хамаарал:} Систем нь зөвхөн интернет сүлжээтэй орчинд бүрэн ажиллах боломжтой.
    \item \textbf{OCR-ийн нарийвчлал:} TensorFlow OCR систем нь одоогоор зөвхөн үндсэн үсгүүдийг таних боломжтой (нарийн бичлэгийн хувилбаруудыг бүрэн дэмжихгүй).
    \item \textbf{Агуулгын хязгаар:} Flashcard dataset нь одоогоор 500+ үгтэй, цаашид өргөжүүлэх шаардлагатай.
    \item \textbf{Хугацааны хязгаар:} Туршилтын үе шат зөвхөн 2--3 сарын хугацаанд хийгдсэн.
\end{itemize}

\subsection{Төгсөлтийн судалгааны ажлын бүтэц}
\label{subsec:structure}

Энэхүү төгсөлтийн судалгааны ажил нь удиртгал, 4 үндсэн бүлэг, дүгнэлт, ном зүй, хавсралт гэсэн хэсгүүдээс бүрдэнэ.

\begin{description}
    \item[Нэгдүгээр бүлэг:] Судалгааны ажлын үндэслэл, зорилго, зорилт, ач холбогдлыг тодорхойлно.
    
    \item[Хоёрдугаар бүлэг:] Монгол бичгийн түүх, flashcard систем, React Native, Firebase, TensorFlow OCR зэрэг онолын судалгаа болон ижил төстэй системүүдийн харьцуулсан шинжилгээг хийнэ.
    
    \item[Гуравдугаар бүлэг:] Системийн шинжилгээ, зохиомж, архитектурын шийдлийг дэлгэрэнгүй танилцуулна.
    
    \item[Дөрөвдүгээр бүлэг:] Системийн хэрэгжүүлэлт, ашигласан технологиуд, кодын тайлбар, туршилтын үр дүнг багтаана.
    
    \item[Дүгнэлт:] Ажлын үр дүнг нэгтгэн дүгнэж, цаашдын хөгжүүлэлтийн саналыг дэвшүүлнэ.
\end{description}

% Хүснэгт: Судалгааны зорилтуудын товч
\begin{table}[htbp]
\centering
\caption{Судалгааны зорилтуудын товч}
\label{tab:research_tasks}
\begin{tabular}{|c|l|l|}
\hline
\textbf{№} & \textbf{Зорилт} & \textbf{Үр дүн} \\
\hline
1 & Онолын судалгаа & Монгол бичиг, flashcard систем судлах \\
\hline
2 & Системийн шинжилгээ & UML диаграммууд, UI/UX загвар \\
\hline
3 & Программ хөгжүүлэлт & React Native аппликейшн \\
\hline
4 & Туршилт ба үнэлгээ & Тестийн үр дүн, хэрэглэгчийн сэтгэл ханамж \\
\hline
\end{tabular}
\end{table}

\vfill
